% Module 12: Industry Applications and Future of Quantum Computing
% Estimated slides: 20 | Duration: ~1.5 hours
% Key topics: Quantum advantage use cases by industry, current state,
%             quantum roadmaps, timeline to fault-tolerant QC,
%             career paths, staying current
% Labs: Research mini-project on a chosen use case

%%%%%%%%%%%%%%%%%%%%%%%%%%%%%%%%%%%%%%%%%%%%%%%%%%%%%%%%%%
\begin{frame}[fragile]\frametitle{}
\begin{center}
{\Large Module 12: Applications and Future of Quantum Computing}
\end{center}
\end{frame}

%%%%%%%%%%%%%%%%%%%%%%%%%%%%%%%%%%%%%%%%%%%%%%%%%%%%%%%%%%%
\begin{frame}[fragile]\frametitle{When Does Quantum Win?}
\begin{itemize}
\item Quantum advantage: a quantum algorithm outperforms the best classical algorithm
\item Three criteria for a good quantum use case:
  \begin{enumerate}
  \item The problem has a quantum algorithm with provable or empirical speedup
  \item The problem size is large enough to benefit from the speedup
  \item The data can be efficiently loaded into the quantum computer
  \end{enumerate}
\item Most real-world problems are still waiting for hardware maturity
\item Near-term advantage: quantum chemistry simulation (50-300 qubits)
\item Medium-term: optimization, finance, machine learning
\end{itemize}
% Suggested diagram: Problem-type vs time-to-advantage matrix for different industries
\end{frame}

%%%%%%%%%%%%%%%%%%%%%%%%%%%%%%%%%%%%%%%%%%%%%%%%%%%%%%%%%%%
\begin{frame}[fragile]\frametitle{Quantum-Inspired Optimization}
\begin{itemize}
\item Do NOT always need full universal quantum computer
\item Specialized optimization hardware can solve useful problems
\item Leverages quantum-inspired algorithms on classical hardware
\item Near-term practical applications available today
\item Bridges gap between current technology and future quantum systems
\item Cost-effective solution for optimization problems
\end{itemize}
\end{frame}

%%%%%%%%%%%%%%%%%%%%%%%%%%%%%%%%%%%%%%%%%%%%%%%%%%%%%%%%%%%
\begin{frame}[fragile]\frametitle{The Ising Machine Concept}
\begin{itemize}
\item Hardware inspired by Ising model from physics
\item Represent optimization as energy minimization
\item System naturally seeks lowest energy state
\item Lowest energy corresponds to best solution
\item Analogy: ball rolling downhill to lowest point
\item Optimization variables evolve toward minimum energy configuration
\item Efficient for combinatorial optimization problems
\end{itemize}
\end{frame}

%%%%%%%%%%%%%%%%%%%%%%%%%%%%%%%%%%%%%%%%%%%%%%%%%%%%%%%%%%%
\begin{frame}[fragile]\frametitle{Hybrid Architecture Workflow}
\begin{itemize}
\item Step 1: Customer defines problem (e.g., delivery routing)
\item Step 2: Translation layer converts to mathematical form (QUBO)
\item Step 3: Specialized hardware searches minimum energy state
\item Step 4: Results returned as optimized solution
\item QUBO: Quadratic Unconstrained Binary Optimization
\item Provides better scheduling and reduced costs
\item Combines classical and quantum-inspired approaches
\end{itemize}
\end{frame}

%%%%%%%%%%%%%%%%%%%%%%%%%%%%%%%%%%%%%%%%%%%%%%%%%%%%%%%%%%%
\begin{frame}[fragile]\frametitle{Performance Benefits}
\begin{itemize}
\item Reported improvements: 8x to 20x speedups
\item Enables solving previously unsolved problems
\item Especially useful for large complex problems
\item Handles problems where classical approaches fail
\item Real-world applications in logistics and finance
\item Demonstrated feasibility for industrial use cases
\end{itemize}
\end{frame}

%%%%%%%%%%%%%%%%%%%%%%%%%%%%%%%%%%%%%%%%%%%%%%%%%%%%%%%%%%%
\begin{frame}[fragile]\frametitle{QuantFluence Product Stack}
\begin{itemize}
\item Three main product areas (representative industry example)
\item A: Quantum Computer (photonic, prototype expected 2029)
\item B: Quantum-Inspired Optimization Machine (near-term monetizable)
\item C: Quantum/Photonic Components (hardware modules)
\item Applications: logistics routing, scheduling, resource allocation
\item Portfolio optimization, manufacturing optimization
\item Revenue generation through component sales
\end{itemize}
\end{frame}

%%%%%%%%%%%%%%%%%%%%%%%%%%%%%%%%%%%%%%%%%%%%%%%%%%%%%%%%%%%
\begin{frame}[fragile]\frametitle{Customer Usage Models}
\begin{itemize}
\item Type 1 Expert: Already understand optimization mapping, use hardware directly
\item Type 2 Assisted: Know problem exists, need help with mathematical conversion
\item Type 3 Full Service: Company handles proof of concept and conversion
\item Partner ecosystem provides support and training
\item Customers gradually build internal capability
\item Flexible engagement models for different maturity levels
\end{itemize}
\end{frame}

%%%%%%%%%%%%%%%%%%%%%%%%%%%%%%%%%%%%%%%%%%%%%%%%%%%%%%%%%%%
\begin{frame}[fragile]\frametitle{Finance: Portfolio Optimization and Risk}
\begin{itemize}
\item Portfolio optimization: choose investments to maximize return, minimize risk
\item NP-hard in general; QAOA provides quantum heuristic
\item Monte Carlo risk assessment: Grover's amplitude estimation gives quadratic speedup
\item Option pricing: quantum speedup in evaluating complex derivatives
\item Early adopters: JPMorgan Chase, Goldman Sachs, BNP Paribas, BBVA
\item Near-term: prototype benchmarks; practical advantage expected post-FTQC
\end{itemize}
% Suggested diagram: Portfolio optimization problem as a graph; quantum vs classical runtime
\end{frame}

%%%%%%%%%%%%%%%%%%%%%%%%%%%%%%%%%%%%%%%%%%%%%%%%%%%%%%%%%%%
\begin{frame}[fragile]\frametitle{Pharmaceuticals and Drug Discovery}
\begin{itemize}
\item Target: simulate molecular interactions at quantum level for drug design
\item FeMo-co simulation (nitrogen fixation): beyond classical, needs $\sim$200 logical qubits
\item Protein-ligand binding affinity: quantum advantage over empirical force fields
\item Timeline: Merck, Pfizer, Roche investing in quantum chemistry partnerships
\item Near-term: VQE demonstrations on small molecules (H$_2$, LiH, BeH$_2$)
\item Long-term (5-10 years): practical molecular simulation at pharma-relevant scale
\end{itemize}
% Suggested diagram: Drug discovery pipeline with quantum simulation step highlighted
\end{frame}

%%%%%%%%%%%%%%%%%%%%%%%%%%%%%%%%%%%%%%%%%%%%%%%%%%%%%%%%%%%
\begin{frame}[fragile]\frametitle{Logistics and Supply Chain}
\begin{itemize}
\item Vehicle routing problem (VRP): assign vehicles to delivery routes efficiently
\item NP-hard: combinatorial explosion with city count
\item QAOA formulation: encode VRP as Max-Cut variant, find approximate optimal routes
\item Quantum annealing (D-Wave): already used by Volkswagen, DHL for logistics
\item Potential: 15-20\% cost reduction in complex logistics networks
\item Near-term: hybrid classical-quantum solvers for small subproblems
\end{itemize}
% Suggested diagram: City-to-city routing problem graph; quantum vs classical solution comparison
\end{frame}

%%%%%%%%%%%%%%%%%%%%%%%%%%%%%%%%%%%%%%%%%%%%%%%%%%%%%%%%%%%
\begin{frame}[fragile]\frametitle{Artificial Intelligence and Machine Learning}
\begin{itemize}
\item Quantum ML: potential speedups in linear algebra, sampling, optimization
\item Quantum principal component analysis: exponential speedup for PCA (if data is in quantum RAM)
\item Quantum support vector machines: quantum kernel trick for feature space classification
\item Quantum neural networks: parameterized circuits as trainable models
\item Dequantization: many claimed quantum ML speedups have been matched classically
\item Practical ML: quantum advantage for ML is still being rigorously established
\item Honest view: most quantum ML advantages are theoretical; hardware not there yet
\end{itemize}
% Suggested diagram: Classical ML pipeline with quantum components highlighted
\end{frame}

%%%%%%%%%%%%%%%%%%%%%%%%%%%%%%%%%%%%%%%%%%%%%%%%%%%%%%%%%%%
\begin{frame}[fragile]\frametitle{Energy and Climate}
\begin{itemize}
\item Battery materials: simulate Li-ion cathode quantum chemistry for better batteries
\item Catalyst design: model nitrogen fixation catalyst to reduce Haber-Bosch energy cost
\item Solar cells: design molecular chromophores for more efficient light harvesting
\item Carbon capture: discover better catalysts for CO$_2$ reduction
\item Grid optimization: power grid scheduling as combinatorial optimization
\item High-temperature superconductors: could revolutionize power transmission
\end{itemize}
% Suggested diagram: Energy applications landscape with quantum simulation touchpoints
\end{frame}

%%%%%%%%%%%%%%%%%%%%%%%%%%%%%%%%%%%%%%%%%%%%%%%%%%%%%%%%%%%
\begin{frame}[fragile]\frametitle{Quantum Computing Roadmaps}
\begin{itemize}
\item \textbf{IBM}: 100k+ physical qubits by 2033; logical qubit experiments by 2029; error-corrected by 2033
\item \textbf{Google}: error below threshold demonstrated (Willow 2024); fault-tolerant timeline TBD
\item \textbf{IonQ}: Forte Pro (2024); Tempe facility for large-scale ion trap systems
\item \textbf{Quantinuum}: System Model H2; 99.99\% 2-qubit fidelity milestone reached
\item \textbf{Microsoft}: topological qubit demonstration; targeting million-qubit roadmap
\item Reality check: timelines have historically been optimistic; Moore-like scaling is hard
\end{itemize}
% Suggested diagram: Timeline chart showing different companies' qubit count vs year roadmaps
\end{frame}

%%%%%%%%%%%%%%%%%%%%%%%%%%%%%%%%%%%%%%%%%%%%%%%%%%%%%%%%%%%
\begin{frame}[fragile]\frametitle{The Path to Fault-Tolerant QC}
\begin{itemize}
\item Stage 1 (now): NISQ devices, 100-1000 noisy qubits, limited depth, no error correction
\item Stage 2 (2025-2030): error mitigation + early error correction; 1000-10k physical qubits
\item Stage 3 (2030-2035): fault-tolerant logical qubits; Shor's possible on cryptographic sizes
\item Stage 4 (2035+): scalable fault-tolerant QC; broad quantum advantage
\item Each stage unlocks new applications: chemistry $\to$ optimization $\to$ ML $\to$ simulation
\item Quantum error correction is the critical unsolved engineering challenge
\end{itemize}
% Suggested diagram: Stage-gate diagram from NISQ to FTQC with unlocked applications
\end{frame}

%%%%%%%%%%%%%%%%%%%%%%%%%%%%%%%%%%%%%%%%%%%%%%%%%%%%%%%%%%%
\begin{frame}[fragile]\frametitle{Career Paths in Quantum Computing}
\begin{itemize}
\item \textbf{Quantum software engineer}: Qiskit, Cirq, PennyLane development; algorithm implementation
\item \textbf{Quantum algorithm researcher}: theoretical; new algorithms, complexity, error correction
\item \textbf{Quantum hardware engineer}: qubit fabrication, control systems, cryogenics
\item \textbf{Quantum application scientist}: map domain problems (chemistry, finance) to quantum algorithms
\item \textbf{Quantum ML researcher}: quantum neural networks, kernel methods, encoding
\item \textbf{Quantum security analyst}: post-quantum cryptography, QKD systems
\item Entry paths: physics, CS/engineering, chemistry; Qiskit certification available
\end{itemize}
% Suggested diagram: Career tree branching from software to hardware to theory to applications
\end{frame}

%%%%%%%%%%%%%%%%%%%%%%%%%%%%%%%%%%%%%%%%%%%%%%%%%%%%%%%%%%%
\begin{frame}[fragile]\frametitle{How to Stay Current}
\begin{itemize}
\item \textbf{arXiv}: quant-ph section; follow quantum computing papers daily
\item \textbf{IBM Quantum Learning}: updated tutorials, labs, and courses (learning.quantum.ibm.com)
\item \textbf{Qiskit blog}: new features, research, industry use cases
\item \textbf{Quantum Computing Stack Exchange}: questions and answers from the community
\item \textbf{Journals}: Nature, PRL, Quantum (open-access)
\item \textbf{Conferences}: QIP, APS March Meeting, IEEE Quantum Week
\item \textbf{Open source}: contribute to Qiskit, Cirq, PennyLane on GitHub
\end{itemize}
\end{frame}

%%%%%%%%%%%%%%%%%%%%%%%%%%%%%%%%%%%%%%%%%%%%%%%%%%%%%%%%%%%
\begin{frame}[fragile]\frametitle{What Are Quantum Computers Good For?}
\begin{itemize}
\item Long term: simulating other quantum systems
\item Pharmaceutical design: molecular simulation
\item Material discovery: new compounds and properties
\item Classical computers struggle with quantum simulation
\item $n$-atom molecule needs roughly $k^n$ amplitudes ($k \geq 2$)
\item Quantum computers: linear scaling with atoms
\item Quantum corollary to Moore's law: exponential advantage
\item Potential for multi-trillion dollar industries over 100+ years
\end{itemize}
\end{frame}


%%%%%%%%%%%%%%%%%%%%%%%%%%%%%%%%%%%%%%%%%%%%%%%%%%%%%%%%%%%
\begin{frame}[fragile]\frametitle{Course Summary: What You Have Learned}
\begin{itemize}
\item \textbf{Foundations}: qubits, superposition, entanglement, measurement, interference, decoherence
\item \textbf{Gates and circuits}: single/multi-qubit gates, circuit diagrams, Bell states
\item \textbf{Algorithms}: Grover's search, Deutsch-Jozsa, QFT, Shor's, VQE, QAOA
\item \textbf{Communication}: teleportation, BB84 QKD, post-quantum cryptography
\item \textbf{Hardware}: superconducting, trapped ion, photonic; IBM Q, Google, Braket
\item \textbf{Programming}: Qiskit circuits, simulators, parameterized circuits, hardware submission
\item \textbf{Applications}: chemistry, finance, logistics, AI; realistic timeline for impact
\end{itemize}
\end{frame}

%%%%%%%%%%%%%%%%%%%%%%%%%%%%%%%%%%%%%%%%%%%%%%%%%%%%%%%%%%%
\begin{frame}[fragile]\frametitle{Mini-Project Ideas}
\begin{itemize}
\item Implement Grover's search for a 4-qubit problem and run on IBM hardware
\item Simulate H$_2$ molecule VQE with Qiskit Nature; vary bond length
\item Implement BB84 protocol simulation with and without eavesdropper
\item Build and benchmark a QAOA solver for a 5-node Max-Cut problem
\item Train a quantum SVM on the Iris dataset using Qiskit ML
\item Compare noise models from different IBM Quantum devices
\end{itemize}
\end{frame}
